\documentclass{article}

\usepackage{amsmath,amssymb}
\usepackage{xurl}
\usepackage[hidelinks]{hyperref}
\setlength{\emergencystretch}{2em}

\input{command}

\title{The Trace Sign in Wolf's Lorentz-Cone Converse}
\author{}
\date{2026-06-23}

\begin{document}
\maketitle
\gapnote{local-correction}{resolved}

\section{The assertion}

Let $A\in\MN{d}$ be Hermitian. Wolf's Proposition~3.7 first states that
positive semidefiniteness implies
\begin{align}
    A\geq 0
    \quad\Longrightarrow\quad
    \tr(A^2)\leq\tr(A)^2.
    \label{eq:wolf_ch3_lorentz_cone_trace_sign_forward_inequality}
\end{align}
The converse printed in the same proposition is
\begin{align}
    \tr(A)^2\geq(d-1)\tr(A^2)
    \quad\Longrightarrow\quad
    A\geq 0.
    \label{eq:wolf_ch3_lorentz_cone_trace_sign_printed_converse}
\end{align}
Both assertions appear in Proposition~3.7 of
Ref.~\cite{Wolf2012Quantum}.\footnote{Local source:
\path{Notes/WolfNoteTexSource/ch03_positive_not_completely.tex},
lines~693--722.}

\section{The point at issue}

The squared trace in
Eq.~\ref{eq:wolf_ch3_lorentz_cone_trace_sign_printed_converse} does not
determine the sign of $\tr(A)$. For $A=-\Id_d$ with $d\geq 1$,
\begin{align}
    A&=-\Id_d\not\geq 0,
    \label{eq:wolf_ch3_lorentz_cone_trace_sign_negative_identity_not_positive}\\
    \tr(A)^2
        &=d^2,
    \label{eq:wolf_ch3_lorentz_cone_trace_sign_negative_identity_trace_square}\\
    (d-1)\tr(A^2)
        &=d(d-1),
    \label{eq:wolf_ch3_lorentz_cone_trace_sign_negative_identity_square_trace}\\
    \tr(A)^2
        &\geq(d-1)\tr(A^2).
    \label{eq:wolf_ch3_lorentz_cone_trace_sign_negative_identity_satisfies_inequality}
\end{align}
Thus $A=-\Id_d$ satisfies the hypothesis of
Eq.~\ref{eq:wolf_ch3_lorentz_cone_trace_sign_printed_converse} but contradicts
its conclusion. The inequality determines a Lorentz cone without selecting
its time orientation; it includes both the positive- and negative-trace
sheets.

The issue also appears in Wolf's proof
\cite[Proposition~3.7]{Wolf2012Quantum}. Write the positive and negative parts
of $A$ as $A_+$ and $A_-$, so that
\begin{align}
    A&=A_+-A_-,
    \label{eq:wolf_ch3_lorentz_cone_trace_sign_positive_negative_decomposition}\\
    \tr(A)
        &=\tr(A_+)-\tr(A_-)
    \label{eq:wolf_ch3_lorentz_cone_trace_sign_trace_decomposition}\\
        &\leq\tr(A_+).
    \label{eq:wolf_ch3_lorentz_cone_trace_sign_trace_upper_bound}
\end{align}
To use the last inequality after squaring, $\tr(A)$ must have the same
orientation as the nonnegative quantity $\tr(A_+)$. The condition
$\tr(A)\geq 0$ supplies that orientation.

\section{Corrected statement}

The valid converse adds a nonnegative-trace hypothesis:
\begin{align}
    A&=A^\dagger,
    &
    \tr(A)&\geq 0,
    &
    (d-1)\tr(A^2)&\leq\tr(A)^2
    \quad\Longrightarrow\quad
    A\geq 0.
    \label{eq:wolf_ch3_lorentz_cone_trace_sign_corrected_converse}
\end{align}
The complex matrix formalization writes the trace conditions using real parts.
For Hermitian $A$, the relevant traces are real, so the formal statement agrees
with Eq.~\ref{eq:wolf_ch3_lorentz_cone_trace_sign_corrected_converse}.\footnote{
Formal declaration:
\leanid{Matrix.IsHermitian.posSemidef_of_trace_re_nonneg_of_card_sub_one_mul_trace_sq_re_le}
in \doclink{QICLean/Analysis/MatrixTraceInequalities}. Tracked by
\href{https://github.com/LionSR/TNLean/issues/3402}{TNLean issue \#3402}.}

The proof reduces the claim to the eigenvalues of $A$. Let
$\lambda_0,\lambda_1,\ldots,\lambda_{d-1}$ be the eigenvalues, suppose
$\lambda_0<0$, and define
\begin{align}
    T&=\sum_{i=0}^{d-1}\lambda_i,
    &
    R&=\sum_{i\neq 0}\lambda_i.
    \label{eq:wolf_ch3_lorentz_cone_trace_sign_eigenvalue_sums}
\end{align}
The trace assumption gives $T\geq 0$, while $T=\lambda_0+R$ and
$\lambda_0<0$ imply
\begin{align}
    0\leq T&<R,
    \label{eq:wolf_ch3_lorentz_cone_trace_sign_strict_sum_comparison}\\
    T^2&<R^2.
    \label{eq:wolf_ch3_lorentz_cone_trace_sign_strict_square_comparison}
\end{align}
Cauchy's inequality for the remaining $d-1$ eigenvalues yields
\begin{align}
    R^2
        &\leq(d-1)\sum_{i\neq 0}\lambda_i^2
    \label{eq:wolf_ch3_lorentz_cone_trace_sign_cauchy_remaining_eigenvalues}\\
        &\leq(d-1)\sum_{i=0}^{d-1}\lambda_i^2.
    \label{eq:wolf_ch3_lorentz_cone_trace_sign_cauchy_all_eigenvalues}
\end{align}
Because $T=\tr(A)$ and $\sum_i\lambda_i^2=\tr(A^2)$, combining
Eqs.~\ref{eq:wolf_ch3_lorentz_cone_trace_sign_strict_square_comparison}
and~\ref{eq:wolf_ch3_lorentz_cone_trace_sign_cauchy_all_eigenvalues} gives
\begin{align}
    \tr(A)^2
        <(d-1)\tr(A^2),
    \label{eq:wolf_ch3_lorentz_cone_trace_sign_contradiction}
\end{align}
contrary to the hypothesis in
Eq.~\ref{eq:wolf_ch3_lorentz_cone_trace_sign_corrected_converse}. Hence no
eigenvalue is negative, and $A\geq 0$.

\section{Conclusion}

The converse in Proposition~3.7 of Ref.~\cite{Wolf2012Quantum} is false
without a trace-sign condition. The local correction is
Eq.~\ref{eq:wolf_ch3_lorentz_cone_trace_sign_corrected_converse}. The formal
development records this trace-nonnegative, future-cone statement for use in
the positive-map discussion and does not mark the unrestricted printed
converse as formalized.

\bibliographystyle{alpha}
\bibliography{references}

\end{document}
